% SOURCES: README.md; draft/Features/features.md; docs/0.2.0v/plan.md (§1 North Star,
%          §2 Scope Tiers); docs/1.0.0v/supervisor-report.md (§1)

\chapter{Introduction}
\label{ch:introduction}

\section{Field of Work}
The modern web is built from three languages --- HTML for structure, CSS for
presentation, and JavaScript for behaviour --- yet authoring a page that is
simultaneously \emph{semantic}, \emph{responsive}, and \emph{accessible} remains a
specialised skill. Over the last decade a class of \emph{visual web builders} ---
tools such as Webflow, Framer, and Wix --- has lowered the barrier to entry by
letting users compose pages on a canvas rather than in a text editor
\cite{webflow,framer,wix}. These tools trade code for direct manipulation, which
is a genuine advance in accessibility of \emph{authoring}; but they frequently do
so at the cost of the \emph{output}: heavy client-side runtimes, non-semantic
``div soup'', absolute positioning that breaks on unexpected viewports, and markup
that is difficult to audit against accessibility standards.

Two forces sharpen the problem. First, the \emph{no-code / low-code} movement has
brought a large, non-technical audience to web authoring, so the quality of what a
builder emits is now the quality that most sites inherit. Second, accessibility
has moved from a courtesy to a requirement: the Web Content Accessibility
Guidelines (WCAG) \cite{wcag21} are referenced by legislation in many
jurisdictions, and semantic HTML \cite{whatwg-html} is the foundation on which
assistive technology depends. In parallel, professional front-end practice has
converged on \emph{design tokens} --- named, themeable primitives for colour,
spacing, and typography --- as the disciplined way to keep a design system
consistent. A tool built today should treat semantics, accessibility, and tokens
as first-class concerns rather than as afterthoughts to be patched on later.

\section{Problem Statement}
The core tension this project addresses is that \emph{visual authoring} and
\emph{high-quality output} are usually treated as opposites. A user who cannot
write code is typically forced to accept whatever markup a builder emits, with no
guarantee of semantic structure, accessibility compliance, or portability, and
often with a hard dependency on the vendor's hosting and runtime. The question the
project sets out to answer is therefore concrete and testable:

\begin{quote}
\itshape Can a canvas-based visual editor produce hand-written-quality output ---
semantic, deterministic, accessibility-gated HTML and CSS with only opt-in
JavaScript --- without requiring the author to write or read code, and without
locking the result to a proprietary platform?
\end{quote}

\section{Proposed Solution}
Draw-to-Web is a desktop application that answers this question by making a single
\emph{Document Model} the only source of truth. The canvas is a rendering of that
model; a deterministic generator walks the same model to emit semantic HTML5,
tokens-driven CSS, and opt-in vetted JavaScript. Every visual property is
token-bindable and per-breakpoint, so theming and responsiveness fall out of the
data model rather than being bolted on. Accessibility is enforced by a hard gate
(\texttt{axe-core}) that blocks export on any critical or serious violation, and
the output is formatted with \texttt{prettier} so that it reads as though a human
wrote it. Because the tool is a local desktop application, the exported bundle is a
plain folder of files with no build step and no mandatory runtime --- fully
portable, with no vendor lock-in.

Version~1.0.0 extends this foundation with three authoring surfaces that make the
tool faster to use without weakening any of its guarantees, because each reuses the
same document operations and export pipeline: \emph{draw-to-create} (draw a
rectangle and the system proposes the element to place), \emph{match-layout}
(sketch a rough structure and adopt a bundled professional design that matches it),
and a headless \emph{Model Context Protocol} (MCP) server that drives the whole
pipeline programmatically \cite{mcp-spec}, so the builder can be scripted by an
agent or a client with no graphical interface.

\section{Report Organisation}
The remainder of this report is organised as follows. Chapter~\ref{ch:objective}
states the primary and specific objectives and the project scope.
Chapter~\ref{ch:methodology} describes the work methodology: the task-driven,
milestone-gated process, the ownership model, and the quality gates.
Chapter~\ref{ch:plan} presents the implementation plan and the deviations from it.
Chapter~\ref{ch:requirements} gathers the functional and non-functional
requirements and the primary use cases. Chapter~\ref{ch:refstudy} reviews similar
projects and positions Draw-to-Web against them. Chapter~\ref{ch:analysis} gives
the system analysis, including the UML class, use-case, and sequence diagrams.
Chapter~\ref{ch:tech} lists the technologies used and justifies each choice.
Chapter~\ref{ch:impl} details the implementation. Chapter~\ref{ch:testing} covers
testing and evaluation against the defined budgets, and Chapter~\ref{ch:results}
reports the results against the objectives, states the limitations honestly, and
concludes.
